% Part: first-order-logic
% Chapter: introduction
% Section: models-theories

\documentclass[../../../include/open-logic-section]{subfiles}

\begin{document}

\olfileid{fol}{int}{mod}

\olsection{مدل‌ها و نظریه‌ها}

پس از تعریف نحو و معناشناسی منطق مرتبهٔ اول، می‌توانیم بررسی خواص
!!{structure}s و مفاهیم معنایی را آغاز کنیم. همچنین می‌توانیم
دستگاه‌های !!{derivation} را تعریف و بررسی کنیم. دربارهٔ مجموعه‌ای از
!!{sentence}s می‌توانیم بپرسیم: کدام !!{structure}s همهٔ !!{sentence}s
آن مجموعه را صادق می‌سازند؟ اگر مجموعه‌ای از !!{sentence}s چون
~$\Gamma$ داده شده باشد، !!a{structure}~$\Struct{M}$ که آن‌ها را ارضا
می‌کند \emph{مدل~$\Gamma$} نامیده می‌شود. ممکن است از~$\Gamma$ آغاز
کنیم و بکوشیم مدل‌هایش را بیابیم---این مدل‌ها چه شکلی دارند؟ چقدر باید
بزرگ یا کوچک باشند؟ اما ممکن است از یک !!{structure} یا مجموعه‌ای از
!!{structure}s آغاز کنیم و بپرسیم: کدام !!{sentence}s در آن‌ها
صادق‌اند؟ آیا !!{sentence}sی وجود دارند که این !!{structure}s را
\emph{مشخصه‌سازی} کنند، به این معنا که دقیقاً در همان‌ها صادق باشند؟
این‌گونه پرسش‌ها قلمرو \emph{نظریهٔ مدل} هستند. آن‌ها همچنین زیربنای
\emph{روش اصل‌موضوعی} را تشکیل می‌دهند: توصیف مجموعه‌ای از
!!{structure}s به‌وسیلهٔ مجموعه‌ای از !!{sentence}s، یعنی اصول موضوع
یک نظریه. این کار با مشاهدهٔ این واقعیت ممکن می‌شود که دقیقاً همان
!!{sentence}sی که در منطق مرتبهٔ اول از اصول موضوع نتیجه می‌شوند، در
همهٔ مدل‌های اصول موضوع صادق‌اند.

به‌عنوان مثالی بسیار ساده، پیش‌ترتیب‌ها را در نظر بگیرید. پیش‌ترتیب
رابطه‌ای چون~$R$ روی مجموعه‌ای چون~$A$ است که هم بازتابی و هم متعدی
باشد. مجموعهٔ~$A$ همراه با رابطه‌ای دوموضعی چون
$R \subseteq A \times A$ روی آن دقیقاً چیزی است که برای ارائهٔ
!!a{structure} در یک زبان مرتبهٔ اول با تنها یک نماد رابطهٔ
دوموضعی~$\Obj P$ نیاز داریم: می‌گذاریم $\Domain{M} = A$ و
$\Assign{\Obj P}{M} = R$. چون $R$ پیش‌ترتیب است، بازتابی و متعدی است،
و می‌توانیم مجموعه‌ای چون~$\Gamma$ از !!{sentence}s منطق مرتبهٔ اول
بیابیم که این مطلب را بیان کنند:
\begin{align*}
  & \lforall[\Obj v_0][\Atom{\Obj P}{\Obj v_0,\Obj v_0}]\\
  & \lforall[\Obj v_0][\lforall[\Obj v_1][\lforall[\Obj v_2][((\Atom{\Obj P}{\Obj v_0,\Obj v_1} \land \Atom{\Obj P}{\Obj v_1,\Obj v_2}) \lif \Atom{\Obj P}{\Obj v_0,\Obj v_2})]]]
\end{align*}
این !!{sentence}s فقط نمادگذاریِ «برای هر~$x$، $Rxx$» ($R$ بازتابی
است) و «هرگاه $Rxy$ و $Ryz$، آنگاه $Rxz$ نیز» ($R$ متعدی است) هستند.
می‌بینیم !!a{structure}~$\Struct{M}$ مدل این دو !!{sentence} در
~$\Gamma$ است هرگاه و تنها هرگاه $R$ (یعنی $\Assign{\Obj P}{M}$)
پیش‌ترتیبی روی~$A$ (یعنی $\Domain{M}$) باشد. به بیان دیگر، مدل‌های
$\Gamma$ دقیقاً پیش‌ترتیب‌ها هستند. هر خاصیت مشترک همهٔ پیش‌ترتیب‌ها که
بتوان آن را در زبان مرتبهٔ اولی با تنها~$\Obj P$ به‌عنوان !!{predicate}
بیان کرد (مانند بازتابی‌بودن و تعدی بالا)، از دو !!{sentence} موجود
در~$\Gamma$ نتیجه می‌شود و برعکس. بنابراین هرچه دربارهٔ مدل‌های
~$\Gamma$ اثبات کنیم، دربارهٔ همهٔ پیش‌ترتیب‌ها اثبات کرده‌ایم.

برای هر نظریه و ردهٔ معینی از مدل‌ها (مانند $\Gamma$ و همهٔ
پیش‌ترتیب‌ها)، پرسش‌های جالبی دربارهٔ آنچه در زبان مرتبهٔ اول متناظر
بیان‌پذیر است و آنچه بیان‌پذیر نیست وجود خواهد داشت. برخی خواص
!!{structure}s برای همهٔ زبان‌ها و رده‌های مدل‌ها جالب‌اند؛ یعنی خواص
مربوط به اندازهٔ !!{domain}. برای نمونه، همیشه می‌توان برای هر
~$n \in \PosInt$ بیان کرد که !!{domain} دقیقاً $n$~!!{element}s دارد.
همچنین با مجموعه‌ای شامل بی‌نهایت !!{sentence}s می‌توان بیان کرد که
!!{domain} نامتناهی است. اما نمی‌توان بیان کرد که دامنه متناهی یا
!!{nonenumerable} است. این نتایج دربارهٔ محدودیت‌های زبان‌های مرتبهٔ
اول پیامد قضیه‌های فشردگی و لوونهایم--اسکولم‌اند.

\end{document}
